\section{The certificate for the threshold}\label{app:threshold}

We prove Theorem~\ref{thm:threshold}. This appendix states the steps, and
\arxivonly{Section}{or:threshold} gives their full proofs. Fix
$a\in(\frac12,1)$ and write $L=\mathcal L_a$ and $\phi(t)=\E e^{itL}$. Put
$\theta_j=1$ for even $j$ and $\theta_j=a$ for odd $j$, let $m_0=m_1=1$ and
$m_n=(na-\theta_n)^{-1}\sum_{j=1}^{n-1}\theta_jm_jm_{n-j}$ for $n\ge2$. Then
$\E L^n=n!\,c_n$ with $c_n=m_n/\Gamma(1+an)$~\cite[Theorem~1.4]{glr2026}. As
$\E e^{t\lvert L\rvert}<\infty$ for real $t$~\cite[Corollary~2.11]{glr2026},
$\phi(t)=\sum_nc_n(it)^n$ for all real $t$. By~\cite[Theorem~1.1 and
(2.10)]{glr2026}, $L$ has the law of $V^aL'+\xi(1-V)^aL''$, with $V$ uniform on
$[0,1]$, $\xi=\pm1$, $L'$ and $L''$ distributed as $L$, and all four
independent. Since $\lvert\phi(-s)\rvert=\lvert\phi(s)\rvert$, this gives
\begin{equation}\label{eq:thr-fixed}
 \lvert\phi(t)\rvert\le\int_0^1\lvert\phi(tv^a)\rvert\,
 \lvert\phi(t(1-v)^a)\rvert\,dv\qquad(t\in\mathbb R).
\end{equation}
The density $g_a$ of $\mathcal L^\circ_a$ has the characteristic function
$\operatorname{Re}\phi$.

\emph{Step 1, one value of $a$.} If $m_j\le AB^j$ for $1\le j\le N$, with
$A=((N+1)a-1)/N$, then $m_n\le AB^n$ for all $n\ge1$, by induction, since
$(na-1)/(n-1)$ increases in $n$. This bounds the tail of the moment series and
of its derivatives by geometric series. Taylor's formula then bounds
$\lvert\phi\rvert$ on cells covering $[0,T_3]$ by $s_j$, from finitely many
coefficients. The next step turns these bounds into a bound for all $t$. Let
$T_3\ge2^aT_2$, with $T_2$ an end point of a cell, and $\beta>3$ with $a\beta>1$.
Let $C$ be the largest $s_jr_j^\beta$ over the cells $[l_j,r_j]$ that meet
$[T_2,T_3]$, and
\[
 R=T_2^{-1/a}\sum_{r_j\le T_2}s_j
 \Bigl(1-\tfrac12(r_j/T_2)^{1/a}\Bigr)^{-a\beta}
 \bigl(r_j^{1/a}-l_j^{1/a}\bigr)+\frac{2^{a\beta}CT_2^{-\beta}}{a\beta-1} .
\]
If $R<1$, then $\lvert\phi(t)\rvert\le Ct^{-\beta}$ for all $t\ge T_3$
(\arxivonly{Lemma}{lem:thr-tail}). The proof puts this bound into the right
side of~\eqref{eq:thr-fixed}, gets back a smaller bound, and closes with a
continuity argument. It is an explicit form of Step~4 in the proof
of~\cite[Theorem~2.5]{glr2026}. By Fourier inversion $g_a$ then has two
continuous derivatives, and $g_a''(0)=-\frac1\pi\int_0^\infty t^2\operatorname{Re}\phi(t)\,dt$. We
integrate the moment series term by term on $[0,T_3]$, and the two error
terms are explicit.

\emph{Step 2, intervals of $a$.} With $a$ as an interval, ball arithmetic
fails, as the series cancels terms of size about $10^{60}$. On
$I=[a_0-w,a_0+w]$ we put $\rho=w/q$ with $0<q<1$ and expand in $a-a_0$. The
coefficients are computed at the exact point $a_0$ by power series arithmetic.
If $h$ is analytic with $\lvert h\rvert\le M$ on the disc
$\lvert z-a_0\rvert\le\rho$, then Cauchy's estimate bounds the remainder after
degree $K$ by $Mq^{K+1}/(1-q)$ on $I$. We bound $M$ by enclosing $c_n(z)$ on a
square that contains the disc, where nothing cancels, and beyond degree $N$ by
the complex form of Step~1 with the product formula for
$\lvert\Gamma(x)/\Gamma(x+iy)\rvert^2$~\cite[5.8.3]{dlmf}. This gives a
polynomial $G$ and a number $\eta$ with $\lvert g_a''(0)-G(a-a_0)\rvert\le\eta$
for every $a\in I$. It also shows that $a\mapsto g_a''(0)$ is continuous on
$I$ (\arxivonly{Proposition}{prop:thr-interval}).

\emph{Step 3, $a$ near $\frac12$.} Put $y=2a-1$ and
$\mu_n=m_ny^{\lfloor n/2\rfloor}$. Then $\mu_2=a$, and for $n\ge3$ the recursion
for $\mu_n$ has the factor $y$ in its terms with $n$ even and $j$ odd, and no
other change. So every $\mu_n$ is analytic for $\operatorname{Re}a>\frac14$
(\arxivonly{Lemma}{lem:thr-scaled}). The variable $X=L/\sigma$ with
$\sigma^2=\E L^2=1/(y\Gamma(2a))$ has
$\E X^n/n!=\mu_n\Gamma(2a)^{\lfloor n/2\rfloor}/\Gamma(1+an)$ for even $n$, and
the same times $\sqrt{y\Gamma(2a)}$ for odd $n$. It satisfies the same
fixed-point equation. Steps~1 and~2 in this variable reach $a=\frac12$.

\emph{Step 4, the Hankel test.} This step uses no Fourier analysis.

\begin{lemma}[a Hankel test]\label{lem:thr-hankel}
Let $s_j=(2j+1)\E L^{2j}$. Suppose that for some $k\ge1$ and
$x\in\mathbb R^{k+1}$ we have $\sum_{i,j\le k}x_ix_js_{i+j}<0$ or
$\sum_{i,j\le k}x_ix_js_{i+j+1}<0$. Then case (a) of
Theorem~\ref{thm:two-modes} fails, so case (b) holds.
\end{lemma}

\begin{proof}
In case (a), $\mathcal L^\circ_a$ is unimodal about $0$. By Khinchine's
theorem~\cite{dj1988} it has the law of $UZ$, with $U$ uniform on
$(0,1)$ and independent of $Z$. So $s_j=(2j+1)\E(\mathcal L^\circ_a)^{2j}
=\E W^j$ with $W=Z^2\ge0$, and the two sums are $\E(\sum_ix_iW^i)^2\ge0$ and
$\E W(\sum_ix_iW^i)^2\ge0$.
\end{proof}

The script \texttt{certify\_threshold.py} checks every hypothesis in ball
arithmetic~\cite{johansson2017}. It evaluates $R$ with $a$ as a ball over $I$
and evaluates $G$ on subintervals of $I$, bisecting until $G\pm\eta$ has one
sign on each. For the Hankel test it takes $k$ and a rational $x$ from a
floating-point eigenvector at the centre of a box of $a$, and checks the sum
on the whole box with a Taylor model as in Step~2. Choices made in floating
point, such as $N$, the precision and $x$, affect only whether a run succeeds.
The cells have half-width $\frac1{32}$, the Taylor order in $t$ is $12$,
$\beta=4$, and the degree $K$ is chosen so that $Mq^{K+1}\le e^{-60}$.
Table~\ref{tab:thr-runs} lists the runs.

\begin{table}[ht]
\centering\small
\begin{tabular}{@{}lllll@{}}
\toprule
range of $a$ & method & pieces & parameters & time\\
\midrule
$[\frac12,0.668]$ & Steps 1 to 3 & 40 & $q=\frac18$, $T_2\in[3,10]$ & 2.1 h\\
$[0.668,0.6712]$ & Steps 1 and 2 & 8 & $q=\frac1{32}$, $\rho\le\frac1{100}$, $T_2\in[4,10]$ & 13 min\\
$[0.6705,1]$ & Step 4 & 115 & $k\le56$ & 25 min\\
\bottomrule
\end{tabular}
\caption{The runs of the certificate. The times are summed over the processes of each run.}\label{tab:thr-runs}
\end{table}

\begin{proof}[Proof of Theorem~\ref{thm:threshold}]
The 40 intervals of the first run cover $[\frac12,0.668]$, and on each the
enclosure of $g_a''(0)$ is negative. The 8 intervals of the second run cover
$[0.668,0.6712]$. On the 2 intervals that cover $[0.668,0.6688]$ the enclosure
is negative, and on the 5 that cover $[0.6692,0.6712]$ it is positive. On
$[0.6688,0.6692]$ we have $\eta<1.1\cdot10^{-15}$. There the enclosure is
negative on closed pieces that end at $z_-$ and positive on closed pieces that
start at $z_+$, and only $(z_-,z_+)$ is undecided. The script prints these
signs and checks this arrangement. So $g_a''(0)<0$ on $(\frac12,z_-]$ and
$g_a''(0)>0$ on $[z_+,0.6712]$. The 115 boxes of the third run cover
$[0.6705,1]$, and on each the first sum of Lemma~\ref{lem:thr-hankel} is
negative.

By Step~1, $g_a$ is even with two continuous derivatives, so $g_a'(0)=0$. If
$g_a''(0)<0$, then $g_a'<0$ on some $(0,\delta)$, which rules out case (b),
where $g_a$ is nondecreasing on $[0,c]$. If $g_a''(0)>0$, then $g_a'>0$ on
some $(0,\delta)$, which rules out case (a). With Lemma~\ref{lem:thr-hankel}
on $[0.6705,1)$ this proves parts (a) and (b). By Step~2,
$a\mapsto g_a''(0)$ is continuous on $[0.6688,0.6692]$, so it has a zero in
$(z_-,z_+)$. The signs above exclude zeros elsewhere in $(\frac12,0.6712]$.
\end{proof}
